\paragraph{Vacuum non-degeneracy and the defect threshold of the displayed potential (conditional; not closure).}
Write \(\mathcal V(\Phi)=\tfrac{\lambda}4\operatorname{Tr}(\Phi^\dagger\Phi-v^2\mathbb I)^2\) for the function of Equation~\eqref{eq:V}, with singular values \(\sigma_i\ge0\), so that \(\mathcal V=\tfrac{\lambda}4\sum_i(\sigma_i^2-v^2)^2\). Two consequences of this form are recorded. (i) The vacuum manifold~\eqref{eq:vacuum-manifold} is \(v\,\mathrm U(4)\): every point has \(\sigma_i=v\), hence \(|\det\Phi|=v^4\ne0\) and \(\Phi\Phi^\dagger=v^2\mathbb I\), so every Repair~II frame~\eqref{eq:vev-is-frame} built on an exact minimum of~\eqref{eq:V} satisfies \(|\det e|=1\), and Assumption~\ref{ass:nondeg} holds automatically on the vacuum manifold of the displayed potential. (ii) The cost of leaving it is sharp: \(\inf\{\mathcal V:\det\Phi=0\}=\lambda v^4/4\), attained exactly at singular values \((0,v,v,v)\), since \(4\mathcal V/\lambda-v^4\) at \(\sigma_1=0\) is a sum of squares; hence at any point where \(\mathcal V<\lambda v^4/4\) the frame~\eqref{eq:vev-is-frame} is non-degenerate, and \(\mathcal V\le\varepsilon<\lambda v^4/4\) forces \(\sigma_i^2\ge v^2-\sqrt{4\varepsilon/\lambda}\) and \(|\det e|\ge(1-\sqrt{4\varepsilon/\lambda}/v^2)^2\). For configurations of Equation~\eqref{eq:Smicro} Assumption~\ref{ass:nondeg} is therefore implied by the pointwise bound \(\mathcal V<\lambda v^4/4\) on the function~\eqref{eq:V}, up to the overall positive measure and trace factors, rather than by an independent non-vanishing hypothesis, and the condensate defects of Chapter~\ref{ch:vierbein-emergence} are events of potential-value cost at least \(\lambda v^4/4\), not flat directions. Along the internal \(\mathrm{SO}(3,1)\) orbit \(\Phi=v\Omega\), \(\mathcal V=(\lambda v^4/4)\operatorname{Tr}(\Omega^{\mathrm T}\Omega-\mathbb I)^2\), equal to \(2\lambda v^4\cosh(2t)\sinh^2t\) for a boost of rapidity \(t\): it vanishes exactly on \(\mathrm{SO}(3)\) and grows without bound along the boosts, so the diagonal vacuum~\eqref{eq:VEV} is a strict global minimum along the non-compact orbit, with quadratic coefficient reproducing the recorded lift of the three boost directions in Section~\ref{sec:status-F}, whose consequences for the Goldstone count are recorded there and are unchanged here. These statements are properties of the displayed Presentation-B potential: that the matrix integral selects it, that the physical state or the bilocal composite of the preceding paragraph realizes the pointwise bound \(\mathcal V<\lambda v^4/4\), and the quantum stability of the minimum are not derived, so Assumptions~\ref{ass:frame-lock} and~\ref{ass:nondeg} are restated as conditional statements and are not discharged, the locking is not derived from Presentation~A, and Problem~D remains OPEN and PROGRAMME-CONDITIONAL.
